% ══════════════════════════════════════════════════════════════
%  Chapter 12 -- Rashomon Sets and Interpretable Trees
% ══════════════════════════════════════════════════════════════

\chapter{Rashomon Sets and Interpretable Trees}\label{ch:rashomon}

You have built a volatility model, run SHAP, and found that VIX is the most
important feature.  You refit on a slightly different training window and now
ATM implied volatility takes over.  A third refit: lagged $\RV_{t-1}$ wins.
Which feature is \emph{actually} important?  The unsettling answer is that any
single model's feature importance is a sample of size one from a large
population of equally-good models.  Drawing conclusions from that single
sample is exactly the kind of reasoning we would reject in any other
statistical context.

This chapter introduces a principled alternative.  Instead of asking ``what
does \emph{my} model say is important?'' we ask ``what do \emph{all}
near-optimal models agree is important?''  The set of all such models is
called the \textbf{Rashomon set}, and the tools for analyzing it give us
something SHAP cannot: provably stable feature importance that does not
change when you perturb the training window by a few days.

\begin{prereq}[What You Need Before This Chapter]
\begin{itemize}[leftmargin=1.8em]
  \item \textbf{Chapter~\ref{ch:trees}}.  Decision tree fundamentals: how
    splits are chosen, how CART grows a tree, what makes a tree
    ``optimal.''  The distinction between greedy (top-down) and globally
    optimal (branch-and-bound) tree construction.
  \item \textbf{SHAP basics}.  The idea that SHAP assigns each feature a
    contribution to each individual prediction, and that global SHAP
    importance is obtained by averaging $|\phi_i|$ across predictions.
    You should know what a SHAP summary plot looks like, but you do
    \emph{not} need to understand the Shapley value derivation in detail.
  \item \textbf{Feature importance concepts}.  Permutation importance
    (shuffle a feature, measure how much loss increases) and split-count
    importance (how often a feature appears in tree splits).  Both are
    covered in Chapter~\ref{ch:trees}.
\end{itemize}
\end{prereq}


% ──────────────────────────────────────────────────────────────
\section{The Problem with Single-Model Explanations}\label{sec:rashomon:problem}
% ──────────────────────────────────────────────────────────────

Feature importance from a single model is not a fact about the data.  It is a
fact about \emph{one model fit to the data}.  If there are many models that
fit the data almost equally well, each may assign very different importances,
and the one we happen to report is determined by arbitrary choices: the random
seed, the exact train/test split boundary, the hyperparameter grid, even the
order in which features are processed.

\subsection{Why SHAP Importance Is Unreliable with Near-Substitute Features}

Consider two features that are highly correlated: VIX (the CBOE volatility
index) and the ATM implied volatility from the SPX options surface.  Both
carry similar information about the market's forward-looking volatility
expectation.  A tree model must choose one or the other at each split.  If the
model happens to split on VIX first, VIX gets credit for the variance it
explains, and ATM IV gets little credit because the residual information it
adds is small.  Reverse the split order and ATM IV gets the credit.

SHAP values \citep{Lundberg2017} partially address this by considering all
feature orderings.  For a model $f$ and a prediction at input $x$, the SHAP
value for feature $i$ is the unique additive attribution satisfying local
accuracy, missingness, and consistency:

\begin{equation}\label{eq:rashomon:shap}
  \phi_i(f, x) = \sum_{S \subseteq F \setminus \{i\}}
    \frac{|S|!\,(|F| - |S| - 1)!}{|F|!}
    \bigl[f_x(S \cup \{i\}) - f_x(S)\bigr],
\end{equation}

\noindent where:
\begin{itemize}[leftmargin=1.8em]
  \item $F$ is the full feature set, $|F| = M$
  \item $S$ is a subset of features \emph{excluding} feature $i$
  \item $f_x(S) = \E[f(z) \mid z_S = x_S]$ is the expected prediction when
    only features in $S$ are observed
\end{itemize}

\begin{intuition}[Why SHAP Doesn't Fully Solve the Problem]
SHAP computes importances for all feature orderings, but it does so for
a \emph{single fixed model} $f$.  The problem is upstream: the model $f$
itself is just one of many near-optimal models.  A different near-optimal
model $f'$ with the same predictive accuracy may produce entirely different
SHAP values, especially for correlated features.  Averaging over orderings
within one model does not average over the space of \emph{models}.
\end{intuition}

\subsection{Importance Instability in Practice}

The instability is not hypothetical.  In volatility forecasting, features
like VIX, $\VVIX$, ATM $\IVol$, the 25-delta skew, and lagged $\RV_{t-1}$
are all moderately to highly correlated.  A gradient-boosted tree fit on
2017--2021 data might rank VIX as the top feature, while the same model fit on
2017--2022 data (just one extra year) might rank $\RV_{t-1}$ first and
demote VIX to third.  The model's accuracy barely changes between these two
fits, but the ``feature importance story'' changes completely.

This creates a serious practical problem.  In an internship setting, you will
present your model's feature importance to a desk or to a portfolio manager.
If the ranking shuffles every time you retrain, your explanation is not
credible.  What you need is a statement like: ``Across \emph{all} models
within 2\% of optimal accuracy, VIX and ATM IV are interchangeable
substitutes, but lagged RV is always important.''  That statement requires
looking at all good models, not just one.


% ──────────────────────────────────────────────────────────────
\section{The Rashomon Set}\label{sec:rashomon:set}
% ──────────────────────────────────────────────────────────────

The idea of many equally-good models was articulated by Leo Breiman, who
called it the \textbf{Rashomon effect}, a reference to the Akutagawa story
(and Kurosawa film) in which multiple witnesses give contradictory but
individually plausible accounts of the same event.  The \textbf{Rashomon set}
formalizes this: it is the collection of all models whose loss is close to
the best achievable loss.

\subsection{Formal Definition}

We need three ingredients: a model class $\mathcal{F}$ (e.g., all decision
trees of depth $\leq d$), a dataset $(\bX, \by)$, and a reference model
$f^*$ that achieves the best loss within $\mathcal{F}$.

\begin{definition}[$\epsilon$-Rashomon Set \citep{XinEtAl2022}]
Given a model class $\mathcal{F}$, a benchmark model $f^* \in \mathcal{F}$,
and a tolerance $\epsilon > 0$, the \textbf{$\epsilon$-Rashomon set} is
\begin{equation}\label{eq:rashomon-set}
  \mathcal{R}(\epsilon, f^*, \mathcal{F})
    = \bigl\{f \in \mathcal{F} \;\big|\;
      L(f) \leq (1 + \epsilon)\,L(f^*)\bigr\},
\end{equation}
where $L(f)$ is the loss (e.g., misclassification rate, squared error) of
model $f$ evaluated on the training data.
\end{definition}

\noindent The symbols:
\begin{itemize}[leftmargin=1.8em]
  \item $\epsilon$: the tolerance parameter.  Setting $\epsilon = 0.02$
    means we accept models within 2\% of the best loss.  Typical values
    in the literature range from $0.01$ to $0.10$ \citep{XinEtAl2022}.
  \item $f^*$: the reference model, usually the empirical risk minimizer
    within $\mathcal{F}$ (e.g., found by GOSDT for optimal sparse trees).
  \item $L(f^*)$: the best achievable loss in the class.
\end{itemize}

For decision trees specifically, \citet{XinEtAl2022} define the objective
function as the misclassification loss plus a sparsity penalty:

\begin{equation}\label{eq:tree-objective}
  \operatorname{Obj}(t, \bX, \by) =
    \underbrace{\frac{1}{n}\sum_{i=1}^{n} \mathbb{1}[\hat{y}_i \neq y_i]}_{\text{misclassification loss}}
    + \underbrace{\lambda \, H_t}_{\text{sparsity penalty}},
\end{equation}

\noindent where:
\begin{itemize}[leftmargin=1.8em]
  \item $\hat{y}_i$ is the prediction of tree $t$ for observation $i$
  \item $H_t$ is the number of leaves in tree $t$
  \item $\lambda > 0$ is a regularization parameter penalizing complexity
\end{itemize}

The Rashomon set threshold then becomes $\theta_\epsilon = (1 + \epsilon)
\times \operatorname{Obj}(t_{\text{ref}}, \bX, \by)$, and any tree $t$ with
$\operatorname{Obj}(t, \bX, \by) \leq \theta_\epsilon$ is in the set
\citep{XinEtAl2022}.

\begin{projectconnection}[Rashomon Sets for Regression]
Our volatility project uses squared-error loss (or QLIKE), not
misclassification.  The Rashomon set concept applies identically: replace
$L(f)$ with $\QLIKE$ or MSE and everything carries through.  The
computational algorithms (TreeFARMS, SPLIT) currently target classification
trees, but the conceptual framework (and the VIC/RID analysis tools) is
loss-function agnostic.
\end{projectconnection}

\subsection{How Large Is the Rashomon Set?}

The hypothesis space of sparse trees is enormous.  For trees of depth at most
4 with only 10 binary features, the number of possible trees exceeds
$9.3 \times 10^{20}$ \citep{XinEtAl2022}.  The Rashomon set, fortunately, is
usually a tiny fraction of this space, but ``tiny fraction'' can still be
a very large number in absolute terms.

\citet{XinEtAl2022} report that on the COMPAS dataset with
$\lambda = 0.005$ and a 15\% Rashomon threshold, the set contains
approximately $10^{12}$ trees.  On smaller datasets (Monk2, Bar), sets of
$10^5$ to $10^8$ trees are typical.  The key insight from their experiments
is that \emph{natural baselines (BART, Random Forest, CART + sampling) find
at best a tiny sliver of the Rashomon set}: they recover only hundreds to
thousands of trees when the true set contains millions or more.

\begin{workedexample}{Two Trees, Different Features, Same Accuracy}
Consider a toy dataset with 100 observations, 3 binary features
($X_1, X_2, X_3$), and a binary label $Y$.  Suppose the optimal tree $t^*$
has 4 leaves, splits on $X_1$ then $X_2$, and achieves
$\operatorname{Obj}(t^*, \bX, \by) = 0.08 + 0.01 \times 4 = 0.12$.

With $\epsilon = 0.05$, the Rashomon threshold is $\theta_{0.05} = 1.05
\times 0.12 = 0.126$.

Now consider tree $t'$ with 3 leaves that splits on $X_3$ first, then $X_1$.
Suppose it achieves $\operatorname{Obj}(t', \bX, \by) = 0.09 + 0.01 \times 3
= 0.12$.

\begin{center}
\begin{tabular}{lcccc}
\toprule
Tree & Splits on & Leaves & Error & Objective \\
\midrule
$t^*$ & $X_1, X_2$ & 4 & $0.08$ & $0.12$ \\
$t'$ & $X_3, X_1$ & 3 & $0.09$ & $0.12$ \\
\bottomrule
\end{tabular}
\end{center}

Both trees are in the Rashomon set.  Tree $t^*$ relies heavily on $X_2$;
tree $t'$ does not use $X_2$ at all but uses $X_3$ instead.  If we report
SHAP importance for $t^*$ alone, we conclude $X_2$ is important.  The
Rashomon set reveals that $X_2$ and $X_3$ are substitutes: the data supports
models using either one.
\end{workedexample}


% ──────────────────────────────────────────────────────────────
\section{Optimal Sparse Decision Trees}\label{sec:rashomon:optimal-trees}
% ──────────────────────────────────────────────────────────────

Before we can enumerate \emph{all} near-optimal trees, we must be able to
find \emph{one} provably optimal tree.

\subsection{Why Not Just Use CART?}

CART \citep{Breiman1984} and other greedy algorithms build trees top-down,
choosing the best split at each node independently.  This is fast
($O(npd)$ for $n$ samples, $p$ features, depth $d$), but greedy splits can
be suboptimal.  \citet{BabbarEtAl2025} show that greedy methods exhibit
an average gap of 1--2 percentage points from the optimum, and on some
datasets (e.g., COMPAS) the gap can reach 10 percentage points.

The problem with greedy construction is that a split that looks best at the
root may set up poor options downstream.  Global optimization avoids this
by searching the \emph{entire} space of trees up to a given depth.

\subsection{Branch-and-Bound with Dynamic Programming}

Modern optimal tree algorithms (GOSDT, MurTree) solve the optimization
problem exactly:

\begin{equation}\label{eq:optimal-tree}
  \mathcal{L}^*(D, d, \lambda) = \min_{T \in \mathcal{T}}
    L(T, D, \lambda) \quad \text{s.t. } \operatorname{depth}(T) \leq d,
\end{equation}

\noindent where $L(T, D, \lambda) = \frac{1}{N}\sum_{i=1}^{N}
\ell(T(\bx_i), y_i) + \lambda\,S(T)$ is the regularized loss and
$S(T)$ is the number of leaves \citep{BabbarEtAl2025}.

The key insight behind branch-and-bound is that the optimal solution for
a dataset $D$ at depth $d'$ depends on the solutions for subsets $D(f)$ and
$D(\bar{f})$ at depth $d'-1$, where $f$ is the splitting feature.  Starting
from the root, the algorithm considers all candidate features, tracking
upper and lower bounds on the objective at each split.  When the lower bound
of a subtree exceeds the current best upper bound, that entire subtree is
pruned.

\subsection{SPLIT: Greedy Where It Doesn't Matter, Optimal Where It Does}

A key empirical observation by \citet{BabbarEtAl2025} is that greedy
splits near the \emph{leaves} are almost always optimal, while greedy splits
near the \emph{root} often deviate from the optimum.  This makes sense: near
the leaves, there are only a few possible splits left, so the greedy choice
among them is unlikely to be far from the best.  Near the root, a bad split
propagates errors through the entire tree.

SPLIT (Sparse Lookahead for Interpretable Trees) exploits this observation.
It takes a \textbf{lookahead depth} parameter $d_l < d$ and performs full
branch-and-bound optimization for splits up to depth $d_l$, then switches to
greedy splitting for the remaining $d - d_l$ levels.

\begin{keyresult}[SPLIT Runtime]
For a dataset with $n$ samples, $k$ binary features, depth budget $d$, and
lookahead depth $d_l$, SPLIT (Algorithm~2 in \citet{BabbarEtAl2025}) has
runtime
\[
  \mathcal{O}\bigl(n(d - d_l)\,k^{d_l + 1} + n\,k^{d - d_l}\bigr).
\]
The polynomial-time variant, LicketySPLIT, achieves $O(nk^2 d^2)$,
comfortably polynomial and dramatically faster than the $\mathcal{O}((2k)^d)$
worst case of fully optimal methods.
\end{keyresult}

\begin{projectconnection}[Interpretable Trees for Vol Forecasting]
For our volatility project, a depth-5 tree with 10--15 features is a
realistic interpretable baseline.  An optimal tree of this size can be found
in under a second with LicketySPLIT.
\end{projectconnection}

\subsection{STreeD: Piecewise-Linear Regression Trees}
\label{sec:rashomon:streed}

The optimal tree methods discussed so far assume each leaf predicts a
\emph{constant}: the mean of the training labels that fall into that leaf.
This creates a staircase-shaped prediction function: the forecast jumps
from one flat value to another at each split boundary.
\citet{VanDenBosEtAl2024} introduce \textbf{STreeD} (STreeD Regression
Trees), a dynamic-programming framework that extends optimal regression
trees to \textbf{piecewise-linear} leaf models.  They develop three methods
of increasing expressiveness:

\begin{enumerate}[leftmargin=1.8em]
  \item \textbf{SRT-C} (piecewise-constant): An improved DP algorithm for
    constant-leaf regression trees with a specialized depth-two solver that
    achieves orders-of-magnitude speedups over previous optimal methods
    (e.g., 18$\times$ faster than OSRT on average).

  \item \textbf{SRT-SL} (simple linear regression): The \emph{first optimal
    method} for piecewise simple linear regression trees.  Each leaf fits a
    one-variable linear model $y = \hat{\beta}_0 + \hat{\beta}_j x_j$,
    selecting the single best feature $j$ for that leaf.  Ridge
    regularization prevents overfitting in small leaves.

  \item \textbf{SRT-L} (multiple linear regression): The \emph{first optimal
    method} for piecewise multiple linear regression trees.  Each leaf fits
    a full linear model with elastic net regularization ($\ell_1 + \ell_2$
    penalty), solved via coordinate descent.
\end{enumerate}

The key idea behind SRT-SL is that each leaf selects \emph{one} continuous
feature and fits a simple linear regression with ridge regularization.  The
leaf-level objective for a leaf containing data subset $\mathcal{D}$ is:

\begin{equation}\label{eq:streed-sl}
  \min_{j,\,\hat{\beta}_0,\,\hat{\beta}_j}
    \sum_{(x,\,b,\,y) \in \mathcal{D}}
    (y - \hat{\beta}_0 - x_j\,\hat{\beta}_j)^2
    + \gamma\,\hat{\beta}_j^2,
\end{equation}

\noindent where:
\begin{itemize}[leftmargin=1.8em]
  \item $j$ is the index of the continuous feature selected for this leaf
  \item $\hat{\beta}_0$ is the intercept,
    $\hat{\beta}_j$ is the slope on feature $x_j$
  \item $\gamma > 0$ is the ridge regularization parameter
  \item The leaf selects whichever feature $j$ yields the smallest
    regularized SSE
\end{itemize}

The closed-form solutions for the optimal coefficients are:
\begin{align}
  \hat{\beta}_j &= \frac{n\sum x_j y - \sum y \sum x_j}
    {n\sum x_j^2 - (\sum x_j)^2 + n\gamma},
    \label{eq:streed-slope}\\
  \hat{\beta}_0 &= \frac{\sum y}{n}
    - \hat{\beta}_j\,\frac{\sum x_j}{n},
    \label{eq:streed-intercept}
\end{align}

\noindent where all sums run over the instances in the leaf and
$n = |\mathcal{D}|$.

\begin{intuition}[Why Piecewise-Linear Leaves Matter]
A constant-leaf tree predicts $\bar{y}$ in each partition, creating a
staircase function.  A linear-leaf tree fits a slope within each
partition, so it captures \emph{local trends}.  For volatility, this means a
leaf can express ``RV increases linearly with VIX within this regime''
rather than just ``RV is high in this regime.''  The tree handles the
nonlinear regime boundaries (splits), and the linear models handle the
smooth within-regime relationships.
\end{intuition}

\textbf{Scalability.}  The depth-two algorithm from
\citet{VanDenBosEtAl2024} is the key to STreeD's performance advantage.
By precomputing per-instance costs (the sums $\sum y$, $\sum y^2$,
$\sum x_j$, $\sum x_j^2$, and $\sum x_j y$ for every feature $j$ and
every data subset defined by pairs of binary splits), the depth-two
solver avoids redundant traversals of the data.  Remarkably, fitting
a simple linear regression model per leaf (SRT-SL) costs almost nothing
extra compared to fitting a constant (SRT-C), because the additional
statistics ($\sum x_j$, $\sum x_j^2$, $\sum x_j y$) can be accumulated
in the same pass.

\textbf{Interpretability.}  Every SRT-SL prediction decomposes into two
transparent components: (1) a root-to-leaf path of binary splits that
determines which regime the input falls into, and (2) a one-variable
linear formula in the leaf that produces the forecast.  This is what the
interpretability literature calls a \textbf{short linear formula}: the
entire model is human-readable and can be written on a single page.

\begin{projectconnection}[STreeD for Volatility Forecasting]
SRT-SL is particularly appealing for our project.  Imagine a depth-3 tree
that splits on lagged-$\RV$ quintile at the root, then on a VIX threshold,
creating 4--8 leaves.  Within each leaf, a simple linear regression on one
feature (e.g., $\RV_{t-1}$ or the VIX level) captures the local
relationship.  The result is an interpretable model that can express
statements like: ``In the high-vol, backwardated-VIX regime, tomorrow's RV
increases by 0.15 for each unit increase in today's RV.''  This is far more
informative than a constant prediction per regime, and every coefficient
has a clear economic interpretation.
\end{projectconnection}


% ──────────────────────────────────────────────────────────────
\section{Enumerating the Rashomon Set}\label{sec:rashomon:enumeration}
% ──────────────────────────────────────────────────────────────

Finding one optimal tree is step one.  The real power comes from finding
\emph{all} near-optimal trees, the complete Rashomon set.  This section
covers three algorithms of increasing scalability.

\subsection{TreeFARMS: Exact Enumeration}\label{sec:rashomon:treefarms}

TreeFARMS (Trees FAst RashoMon Sets) by \citet{XinEtAl2022} was the first
algorithm to \emph{completely} enumerate the Rashomon set for sparse decision
trees.  It builds on the GOSDT branch-and-bound framework and modifies it in
two key ways:

\begin{enumerate}[leftmargin=1.8em]
  \item \textbf{Rashomon pruning.}  Instead of pruning subproblems whose
    lower bound exceeds the optimal objective (as in standard GOSDT),
    TreeFARMS prunes those whose lower bound exceeds the \emph{Rashomon
    threshold} $\theta_\epsilon$.  This retains more of the search space:
    exactly the near-optimal region.

  \item \textbf{Return all models.}  Instead of returning only the single
    best tree, TreeFARMS returns every tree in the Rashomon set, stored in a
    compact \textbf{Model Set} data structure.
\end{enumerate}

The Model Set representation is critical for scalability.  A
\textbf{Model Set Instance (MSI)} represents a subproblem paired with its
objective value.  Many trees share identical subtrees, and the Model Set
exploits this by storing shared components once.  The loss function for
decision trees takes on a discrete set of values (approximately $n$ distinct
values for $n$ training samples), while the number of trees in the Rashomon
set can be orders of magnitude larger.  By grouping trees with the same
objective, TreeFARMS avoids massive data duplication.

\begin{keyidea}[TreeFARMS: From Optimization to Feasibility]
Standard ML algorithms solve an \emph{optimization} problem: find the single
best model.  TreeFARMS solves a \emph{feasibility} problem: find all models
within $\epsilon$ of the best.  This reframing is the core contribution.
Perhaps the tiny sacrifice in empirical risk makes the difference between a
model that can be used (interpretable, fair, aligned with domain knowledge)
and one that cannot.
\end{keyidea}

\textbf{Limitations.}  TreeFARMS provides exact enumeration but its runtime
and memory scale exponentially with tree depth and the number of features.
On the Bike dataset ($n \approx 17{,}000$, $k = 60$ binary features, depth
5), TreeFARMS requires approximately 700 seconds and over 50~GB of memory
\citep{HeileBabbar2025}.  On larger datasets, it runs out of memory entirely.

\subsection{RESPLIT: Fast Approximation via Greedy Leaves}
\label{sec:rashomon:resplit}

\citet{BabbarEtAl2025} extend SPLIT to Rashomon set computation with the
RESPLIT algorithm.  The idea is simple: use SPLIT to find a set of
``prefix trees'' (partial trees optimized to the lookahead depth), then call
TreeFARMS on each prefix to enumerate the near-optimal completions.

Because SPLIT uses greedy splits near the leaves, RESPLIT does not
exhaustively search the full tree space.  This makes it an
\emph{approximation}: it may miss some trees in the true Rashomon set.
However, \citet{BabbarEtAl2025} show that the approximation is remarkably
accurate.  On six benchmark datasets, the Pearson correlation between
variable importances computed from the RESPLIT-approximated Rashomon set
and the full Rashomon set is nearly 1.0:

\begin{center}
\begin{tabular}{lrrrr}
\toprule
Dataset & Full (s) & RESPLIT (s) & Speedup & $\tau$ (correlation) \\
\midrule
COMPAS       & 152    & 18   & $8\times$   & 1.000 \\
Spambase     & 2{,}659 & 154  & $17\times$  & 0.930 \\
Netherlands  & 4{,}255 & 216  & $20\times$  & 0.932 \\
HELOC        & 5{,}564 & 337  & $17\times$  & 0.979 \\
HIV          & 9{,}273 & 388  & $24\times$  & 0.959 \\
Bike         & 14{,}330 & 194 & $74\times$  & 0.999 \\
\bottomrule
\end{tabular}
\end{center}
\emph{Source: Table~1 in \citet{BabbarEtAl2025}.  Parameters:
10 bootstrapped datasets, $\lambda = 0.02$, $\epsilon = 0.01$, depth 5,
lookahead depth 2.}

\subsection{LicketyRESPLIT: Polynomial-Time Approximation}
\label{sec:rashomon:licketyresplit}

\citet{HeileBabbar2025} push the scalability further with LicketyRESPLIT,
which replaces TreeFARMS entirely with a polynomial-time enumeration
strategy.  At each node, the algorithm considers all features and prunes
splits whose LicketySPLIT-completed cost would exceed the Rashomon budget
$B = (1 + \epsilon) \cdot \text{LicketySPLIT}(D, \lambda, d)$.  It then
recursively enumerates left and right subtrees, filtering by the remaining
budget.

The key theoretical results are:

\begin{keyresult}[LicketyRESPLIT Complexity \citep{HeileBabbar2025}]
Given a dataset $D$ of size $n$ with $k$ features and max depth $d$:
\begin{itemize}[leftmargin=1.8em]
  \item \textbf{Runtime:} $\mathcal{O}(|R|\,n\,k^3\,d^3)$, where $|R|$ is
    the size of the recovered Rashomon set.  This is polynomial in $n$, $k$,
    and $d$, and linear in $|R|$.
  \item \textbf{Memory:} $\mathcal{O}\bigl(nk + \sum_{f \in R} S(f)\bigr)$,
    proportional to the input size plus the output size.
\end{itemize}
\end{keyresult}

In practice, LicketyRESPLIT achieves order-of-magnitude improvements over
both TreeFARMS and RESPLIT:

\begin{center}
\begin{tabular}{lrrr}
\toprule
Dataset & \multicolumn{1}{c}{LicketyRESPLIT} & \multicolumn{1}{c}{TreeFARMS} & \multicolumn{1}{c}{RESPLIT} \\
        & Time / RAM & Time / RAM & Time / RAM \\
\midrule
Bike       & \textbf{18.8\,s / 438\,MB} & 685\,s / 51\,GB & 184\,s / 528\,MB \\
Bank       & \textbf{123\,s / 776\,MB}   & OOM              & 238\,s / 2\,GB \\
Covertype  & \textbf{507\,s / 1.8\,GB}   & 1819\,s / 68\,GB & 1295\,s / 3\,GB \\
Student    & \textbf{1.7\,s / 370\,MB}   & 351\,s / 4.7\,GB & 4.0\,s / 382\,MB \\
\bottomrule
\end{tabular}
\end{center}
\emph{Source: Table~1 in \citet{HeileBabbar2025}.  Parameters: $\lambda = 0.01$,
$\epsilon_{\text{mult}} = 0.01$, max depth = 5.}

Despite being an approximation, LicketyRESPLIT achieves near-perfect
precision and recall relative to the true Rashomon set.  On six benchmark
datasets, precision is $\geq 0.91$ and recall is $\geq 0.90$ across all
tested configurations \citep{HeileBabbar2025}.

\begin{warning}[Exact vs.\ Approximate Enumeration]
TreeFARMS gives you the \emph{exact} Rashomon set: every tree in the set, and
no tree outside it.  RESPLIT and LicketyRESPLIT are approximations: they
may miss a few trees (recall $< 1$) and occasionally include a tree slightly
outside the boundary (precision $< 1$).  For downstream variable importance
analysis, this distinction rarely matters: the approximate sets produce
virtually identical importance rankings.  But if you need a formal guarantee
(``no tree in the Rashomon set uses feature $X_7$''), only exact enumeration
suffices.
\end{warning}

Rashomon set algorithms are developing rapidly.
Figure~\ref{fig:rashomon-algorithms} summarizes the three approaches and
their trade-offs.

\begin{figure}[ht]
\centering
\begin{tikzpicture}[
  algobox/.style={draw, rounded corners, minimum width=3.2cm,
    minimum height=1.6cm, align=center, font=\small},
  >=Stealth,
]
% TreeFARMS
\node[algobox, fill=defblue!10] (tf) at (0, 0) {
  \textbf{TreeFARMS}\\[2pt]
  {\scriptsize Exact enumeration}\\
  {\scriptsize Exponential in $d, k$}
};
% RESPLIT
\node[algobox, fill=keyorange!10] (rs) at (5.5, 0) {
  \textbf{RESPLIT}\\[2pt]
  {\scriptsize Greedy leaves + TreeFARMS}\\
  {\scriptsize $10$--$20\times$ faster}
};
% LicketyRESPLIT
\node[algobox, fill=intgreen!10] (lr) at (11, 0) {
  \textbf{LicketyRESPLIT}\\[2pt]
  {\scriptsize Polynomial-time}\\
  {\scriptsize $100\times$ less memory}
};

% Arrows
\draw[->, thick] (tf) -- (rs) node[midway, above, font=\scriptsize] {approximate};
\draw[->, thick] (rs) -- (lr) node[midway, above, font=\scriptsize] {remove TreeFARMS};

% Labels below
\node[font=\scriptsize, text=defblue] at (0, -1.3) {NeurIPS 2022};
\node[font=\scriptsize, text=keyorange] at (5.5, -1.3) {ICML 2025};
\node[font=\scriptsize, text=intgreen] at (11, -1.3) {NeurIPS 2025 Workshop};
\end{tikzpicture}
\caption{Evolution of Rashomon set enumeration algorithms for sparse
decision trees.  Each generation trades a small amount of exactness for
dramatic improvements in scalability.}
\label{fig:rashomon-algorithms}
\end{figure}


% ──────────────────────────────────────────────────────────────
\section{What the Rashomon Set Reveals}\label{sec:rashomon:analysis}
% ──────────────────────────────────────────────────────────────

Two complementary tools turn the raw set of models into concrete
conclusions about feature importance.

\subsection{Model Class Reliance (MCR)}\label{sec:rashomon:mcr}

The simplest analysis is to ask: for each feature, what is the \emph{range}
of its importance across all models in the Rashomon set?  This is
\textbf{Model Class Reliance (MCR)}, introduced by \citet{DongRudin2020}
(building on the concept from Fisher et al., 2019).

For a single model $f$, the \textbf{model reliance} of feature $j$ is
defined as the ratio of the loss when feature $j$ is randomly permuted to
the original loss:

\begin{equation}\label{eq:model-reliance}
  mr_j^{\text{ratio}}(f) = \frac{L(f;\,[X_{\setminus j},\, \bar{X}_j],\, Y)}
    {L(f;\, X,\, Y)},
\end{equation}

\noindent where:
\begin{itemize}[leftmargin=1.8em]
  \item $X_{\setminus j}$ denotes all features except feature $j$
  \item $\bar{X}_j$ is an independent copy of $X_j$ (i.e., feature $j$ is
    reshuffled)
  \item A reliance of 1.0 means the feature contributes nothing; larger
    values indicate greater importance
\end{itemize}

MCR defines two key quantities for each feature $j$:
\begin{align}
  \text{MCR}_{-}(j) &= \min_{f \in \mathcal{R}} mr_j(f), \label{eq:mcr-min}\\
  \text{MCR}_{+}(j) &= \max_{f \in \mathcal{R}} mr_j(f). \label{eq:mcr-max}
\end{align}

The interpretation is clean:
\begin{itemize}[leftmargin=1.8em]
  \item A feature with large $\text{MCR}_{-}$ is important in \emph{every}
    well-performing model.  It is robustly important.
  \item A feature with small $\text{MCR}_{+}$ is unimportant in every
    well-performing model.  It can be safely excluded.
  \item A feature with small $\text{MCR}_{-}$ but large $\text{MCR}_{+}$ is
    a \emph{substitute}: important in some models, not in others.  It can
    be swapped for a correlated alternative without loss of accuracy.
\end{itemize}

\citet{XinEtAl2022} showed that TreeFARMS enables \emph{exact} MCR
computation for decision trees by directly calculating variable importance
for every tree in the set and then finding the min and max.

\subsection{Variable Importance Clouds (VIC)}\label{sec:rashomon:vic}

MCR gives a one-dimensional summary (min and max importance) for each
feature independently.  \textbf{Variable Importance Clouds (VIC)}
\citep{DongRudin2020} capture the full joint structure.

\begin{definition}[Variable Importance Cloud \citep{DongRudin2020}]
The \textbf{model reliance function} $MR : \mathcal{F} \to \R^p$ maps each
model to a vector of its reliances on all $p$ features:
\[
  MR(f) = \bigl(mr_1(f),\; mr_2(f),\; \ldots,\; mr_p(f)\bigr).
\]
The \textbf{Variable Importance Cloud} of the Rashomon set $\mathcal{R}$ is
the set of all such vectors:
\begin{equation}\label{eq:vic}
  \operatorname{VIC}(\mathcal{R}) = \bigl\{MR(f) : f \in \mathcal{R}\bigr\}.
\end{equation}
\end{definition}

The VIC lives in $p$-dimensional space, where each axis represents the
importance of one feature.  To visualize it, \citet{DongRudin2020} project
the VIC onto all pairs of features, producing \textbf{Variable Importance
Diagrams (VIDs)}, 2D scatter plots that reveal substitution patterns:

\begin{itemize}[leftmargin=1.8em]
  \item \textbf{Non-overlapping projections} along one axis: the two
    features have robustly distinct importance levels.  One is always more
    important than the other, regardless of model choice.
  \item \textbf{Overlapping projections}: the features are substitutes.
    Some near-optimal models rely on one, some on the other.
  \item \textbf{Negative slope} in the 2D projection: the features are
    direct substitutes; as one's importance increases, the other's
    decreases.  This is the signature of correlated features competing for
    the same splits.
\end{itemize}

\begin{workedexample}{VIX, VVIX, and ATM IV as Near-Substitutes}
Imagine constructing the Rashomon set for a volatility classification task
(``will tomorrow's RV exceed its 90th percentile?'') using depth-4 trees
with the following features: lagged $\RV_{t-1}$, $\RV_{t-5}$, VIX, $\VVIX$,
and ATM $\IVol$.

After enumeration, we compute $MR(f)$ for each tree $f$ in the set and
project onto two pairs:

\textbf{Pair 1: VIX vs.\ ATM IV.}  The 2D cloud shows a strong negative
slope: when a tree relies heavily on VIX, it does not rely on ATM IV, and
vice versa.  The projections onto each axis overlap substantially.  These
features are \emph{substitutes}: the data supports using either one, but not
both simultaneously.

\textbf{Pair 2: Lagged $\RV_{t-1}$ vs.\ VIX.}  The cloud is a compact blob
with no clear slope.  The projection of $\RV_{t-1}$ onto its axis is
consistently high (all models rely on it), while VIX varies.  Lagged RV is
\emph{robustly important}; VIX is a useful but substitutable addition.

\medskip
\noindent\textbf{Conclusion for feature selection:}  We must include
$\RV_{t-1}$.  We may include either VIX or ATM IV (but including both adds
redundancy, not information).  $\VVIX$ may or may not add value depending on
what else is in the model.
\end{workedexample}

\begin{keyidea}[VIC vs.\ Bootstrapped SHAP]
A natural reaction is: ``I could just bootstrap the data, refit SHAP each
time, and get a distribution of importances.''  This is fundamentally
different from VIC.  Bootstrapped SHAP resamples the \emph{data} but always
uses the same model class and always reports the importance of the single
best model within that class.  VIC holds the data fixed and varies the
\emph{model}: it considers every near-optimal model, not just the best one.

The distinction matters because bootstrap variability conflates data
uncertainty with model uncertainty.  VIC isolates model uncertainty: ``given
this exact dataset, how much does the importance depend on which near-optimal
model we happened to pick?''  This is the right question when your concern
is the stability of your explanations, not the variability of your data.
\end{keyidea}

\subsection{Rashomon Importance Distributions (RID)}
\label{sec:rashomon:rid}

MCR and VIC answer the question ``what is the range of importances across
all near-optimal models \emph{for this dataset}?''  But this answer is
fragile: the Rashomon set itself changes when the dataset is perturbed.
\citet{DonnellyEtAl2023} show that both the Rashomon set size and the MCR
range can vary wildly across bootstrap resamples of the \emph{same} data,
making MCR and VIC unstable summaries of feature importance.

The \textbf{Rashomon Importance Distribution (RID)} proposed by
\citet{DonnellyEtAl2023} addresses this by combining bootstrap resampling
\emph{with} Rashomon set analysis in a two-level procedure that captures
both sources of uncertainty simultaneously.

\textbf{The five-step RID pipeline} (illustrated in Figure~2 of
\citealt{DonnellyEtAl2023}):

\begin{enumerate}[leftmargin=1.8em]
  \item \textbf{Bootstrap.}  Draw $B$ bootstrap datasets
    $\mathcal{D}_1^{(n)}, \ldots, \mathcal{D}_B^{(n)}$ from the original
    data $\mathcal{D}^{(n)}$ (sampling $n$ observations with replacement).
  \item \textbf{Find Rashomon set.}  For each bootstrap dataset
    $\mathcal{D}_b^{(n)}$, compute the Rashomon set
    $\mathcal{R}_{\mathcal{D}_b}^{\varepsilon}$: all models in
    $\mathcal{F}$ whose loss is within $\varepsilon$ of the best model
    \emph{on that bootstrap sample}.
  \item \textbf{Find importances.}  For each model $f$ in each Rashomon
    set, compute the variable importance $\phi_j(f, \mathcal{D}_b^{(n)})$
    using any importance metric (permutation importance, SHAP, model
    reliance, etc.).
  \item \textbf{Find CDF.}  For each bootstrap $b$, compute the
    empirical CDF of importances across the Rashomon set for that bootstrap.
  \item \textbf{Find PDF.}  Average the CDFs across all $B$ bootstraps to
    obtain the RID, then differentiate to get the marginal density.
\end{enumerate}

Formally, RID is defined as a cumulative distribution function.  For
feature $j$, the RID at threshold $k$ measures the expected fraction of
near-optimal models whose importance for feature $j$ is at most $k$,
where the expectation is over the data distribution:

\begin{equation}\label{eq:rid-formal}
  \text{RID}_j(k)
    = \E_{\mathcal{D}_b^{(n)} \sim \mathcal{P}_n}\!\left[
      \frac{|\{f \in \mathcal{R}_{\mathcal{D}_b}^{\varepsilon}
        : \phi_j(f, \mathcal{D}_b^{(n)}) \leq k\}|}
        {|\mathcal{R}_{\mathcal{D}_b}^{\varepsilon}|}
    \right],
\end{equation}

\noindent where:
\begin{itemize}[leftmargin=1.8em]
  \item $\mathcal{P}_n$ is the empirical distribution of the data (i.e.,
    the bootstrap distribution)
  \item $\mathcal{R}_{\mathcal{D}_b}^{\varepsilon}$ is the Rashomon set
    for bootstrap sample $\mathcal{D}_b^{(n)}$
  \item $\phi_j(f, \mathcal{D}_b^{(n)})$ is the importance of feature $j$
    for model $f$ evaluated on $\mathcal{D}_b^{(n)}$
  \item $k$ ranges over $[\phi_{\min}, \phi_{\max}]$, the support of the
    importance metric
\end{itemize}

In practice, we estimate RID by replacing the expectation with a sample
average over $B$ bootstrap replicates:

\begin{equation}\label{eq:rid-empirical}
  \widehat{\text{RID}}_j(k)
    = \frac{1}{B}\sum_{b=1}^{B}
      \frac{|\{f \in \mathcal{R}_{\mathcal{D}_b}^{\varepsilon}
        : \phi_j(f, \mathcal{D}_b^{(n)}) \leq k\}|}
        {|\mathcal{R}_{\mathcal{D}_b}^{\varepsilon}|}.
\end{equation}

\begin{intuition}[In Plain English]
RID asks: ``If I drew a new dataset from the same population, computed all
near-optimal models on it, and measured feature $j$'s importance in each of
those models, what distribution of importances would I see?''  The bootstrap
handles the ``new dataset'' part; the Rashomon set handles the ``all
near-optimal models'' part.  The result is a CDF that captures \emph{both}
data uncertainty and model-choice uncertainty in a single object.
\end{intuition}

\begin{warning}[Additive vs.\ Multiplicative Rashomon Threshold]
\citet{DonnellyEtAl2023} define the Rashomon set using an
\textbf{additive} threshold: $\mathcal{R}^{\varepsilon} = \{f \in
\mathcal{F} : \ell(f) \leq \min_{f'}\ell(f') + \varepsilon\}$.  This
contrasts with the \textbf{multiplicative} threshold used by
\citet{XinEtAl2022} earlier in this chapter: $L(f) \leq (1 + \epsilon)
L(f^*)$.  The two formulations are equivalent when rescaled appropriately,
but be careful not to mix them: an additive $\varepsilon = 0.05$ and a
multiplicative $\epsilon = 0.05$ define different sets.
\end{warning}

\textbf{Why not just bootstrap?}  A natural alternative is naive bootstrap
importance: draw $B$ bootstrap samples, fit the best model on each, and
collect the importances.  This captures data uncertainty but uses only
\emph{one} model per resample (the best one), ignoring the many
near-optimal alternatives.  RID considers \emph{all} near-optimal models
per resample, capturing model-choice uncertainty that naive bootstrapping
misses entirely.

\textbf{Why not just MCR/VIC?}  MCR and VIC analyze the Rashomon set of a
\emph{single} dataset.  They capture model-choice uncertainty but are
blind to data uncertainty.  \citet{DonnellyEtAl2023} demonstrate that MCR
ranges are unstable across bootstrap resamples: for one variable on the
Monk~3 dataset, the MCR range is $[-0.1, 0.33]$ on one resample and
$[0.33, 0.36]$ on another: contradictory conclusions from the same
underlying data.

\begin{keyidea}[RID = Bootstrap $\times$ Rashomon]
RID is the only method that captures \emph{both} sources of uncertainty
simultaneously:
\begin{center}
\begin{tabular}{lccc}
\toprule
\textbf{Method} & \textbf{Data uncertainty} & \textbf{Model uncertainty}
  & \textbf{Stable?} \\
\midrule
Single-model SHAP & No  & No  & No \\
Bootstrap importance & Yes & No  & Partially \\
MCR / VIC            & No  & Yes & No \\
\textbf{RID}         & \textbf{Yes} & \textbf{Yes} & \textbf{Yes} \\
\bottomrule
\end{tabular}
\end{center}
\end{keyidea}

\textbf{Finite-sample guarantees.}  \citet{DonnellyEtAl2023} prove
(Theorem~2) that the estimated $\widehat{\text{RID}}_j$ converges to the
true $\text{RID}_j$ at a rate controlled by the number of bootstraps:
with probability at least $1 - \delta$,
\[
  \bigl|\widehat{\text{RID}}_j(k) - \text{RID}_j(k)\bigr| \leq t
  \quad \text{for all } k,
  \quad \text{when } B \geq \frac{1}{2t^2}\ln\frac{2}{\delta}.
\]
For example, $B = 471$ bootstraps guarantees that the estimated RID is
within $t = 0.075$ of the true value with 90\% confidence.  This is a
practical guarantee: a few hundred bootstraps suffice for reliable results.

\textbf{Metric-agnostic.}  RID works with any variable importance metric
$\phi_j$: permutation importance, SHAP, model reliance, conditional
model reliance, or any other metric with a bounded range.  The framework
treats $\phi_j$ as a black box, making it immediately compatible with
existing importance tools.

\textbf{Empirical stability.}  On four synthetic data-generating
processes, \citet{DonnellyEtAl2023} find that RID achieves a median
Jaccard similarity of 0.69 across independently generated datasets,
compared to below 0.55 for both MCR and VIC.  RID is the only method
whose importance intervals consistently overlap across independent
datasets drawn from the same distribution.

\begin{projectconnection}[RID for Volatility Feature Selection]
For our project, RID would answer questions that neither SHAP nor MCR can:
``Across bootstrap resamples of the training data \emph{and} across all
near-optimal models within each resample, what is the distribution of
VIX's importance?''  If the RID CDF for VIX rises steeply near zero, VIX
is unimportant in most models across most resamples, and we can drop it.
If the CDF rises steeply above 0.2, VIX is robustly important.  If the
CDF rises gradually from 0 to 0.4, VIX is a substitute: sometimes
important, sometimes not, depending on both the data sample and the model
chosen.  Unlike MCR, this conclusion is \emph{stable}: it will not
flip if we add one more month of data.
\end{projectconnection}


% ──────────────────────────────────────────────────────────────
\section{Rashomon Analysis for Volatility Forecasting}
\label{sec:rashomon:vol}
% ──────────────────────────────────────────────────────────────

No published work has applied Rashomon set analysis to financial time-series
forecasting.  This is an open frontier, and it is directly relevant to our
project.  This section describes how the tools from
Sections~\ref{sec:rashomon:enumeration}--\ref{sec:rashomon:analysis} could be
deployed for realized volatility forecasting.

\subsection{Regime-Stable Feature Selection}

Volatility features that matter in a low-vol regime (2017--2019) may not
matter in a crisis regime (March 2020) or a post-crisis recovery (2021).
Single-model SHAP run on the full sample averages over these regimes,
potentially concluding that a feature is ``moderately important'' when it is
actually critical in crises and irrelevant otherwise.

A Rashomon-based approach to regime-stable feature selection:

\begin{enumerate}[leftmargin=1.8em]
  \item \textbf{Rolling-window Rashomon sets.}  For each rolling window
    (e.g., 252-day training sets advanced by one month), compute the
    Rashomon set and the MCR for each feature.
  \item \textbf{Intersection across regimes.}  A feature that has
    $\text{MCR}_{-} > 1$ (i.e., minimum importance above the baseline) in
    \emph{every} window is regime-stable.  A feature whose MCR range
    includes 1.0 in some windows is regime-dependent.
  \item \textbf{Feature selection rule.}  Include only regime-stable
    features in the final model.  For regime-dependent features, consider
    regime-conditional inclusion (e.g., add VIX term structure only when
    VIX $> 20$).
\end{enumerate}

\subsection{Prediction Multiplicity: Quantifying Model-Choice Uncertainty}

The Rashomon set also reveals \textbf{prediction multiplicity}: the range of
forecasts that different near-optimal models produce for the same input.  For
a given feature vector $\bx_t$ (today's features), we can compute
$\{f(\bx_t) : f \in \mathcal{R}\}$ and report the min, max, and spread.

If all near-optimal models agree that tomorrow's volatility will be high,
that is a strong signal.  If some predict high and others predict low, the
forecast is sensitive to model choice, and we should report wider
confidence intervals or flag the day as uncertain.

This is conceptually similar to the disagreement among models in an ensemble,
but with one important difference: ensemble members are not guaranteed to be
near-optimal, while Rashomon set members are.

\subsection{Defensible Presentations}

In a Goldman Sachs internship, you will present model results to people who
will ask hard questions: ``Why did your model use VIX and not ATM IV?'' or
``Your SHAP plot from last month looked completely different.''

Rashomon analysis gives you defensible answers:

\begin{itemize}[leftmargin=1.8em]
  \item ``I examined all 14{,}000 decision trees within 2\% of optimal
    accuracy.  Lagged RV is the top feature in every single one of them.
    VIX and ATM IV are interchangeable; the data supports either, so I
    chose VIX for simplicity.''
  \item ``The feature importance is \emph{provably stable}: no
    near-optimal model exists that does not rely on lagged RV.''
  \item ``The prediction range across all near-optimal models for tomorrow is
    $[0.00012, 0.00018]$.  The spread tells you how much of the forecast
    uncertainty comes from model choice rather than data noise.''
\end{itemize}


% ──────────────────────────────────────────────────────────────
\section{Summary and Connections}\label{sec:rashomon:summary}
% ──────────────────────────────────────────────────────────────

This chapter introduced a fundamentally different approach to model
interpretability.  Rather than explaining one model after the fact (SHAP), we
examine the entire space of near-optimal models before committing to one.

\begin{center}
\begin{tabular}{p{4.2cm}p{4.8cm}p{4.8cm}}
\toprule
\textbf{Aspect} & \textbf{Single-Model (SHAP)} & \textbf{Rashomon Set} \\
\midrule
What is explained & One model's predictions & All near-optimal models \\
Feature importance & Point estimate for one model & Range $[\text{MCR}_{-}, \text{MCR}_{+}]$ \\
Stability & Changes with refit & Provably stable within $\epsilon$ \\
Substitute detection & Not possible & Overlapping VIC projections \\
Model-choice uncertainty & Not quantified & Prediction multiplicity \\
Computational cost & Fast (one model) & Hours (enumeration needed) \\
\bottomrule
\end{tabular}
\end{center}

\textbf{Looking ahead.}  Chapter~\ref{ch:deeplearning} will introduce deep
learning methods for volatility, which are powerful but even harder to
interpret than tree ensembles.  The Rashomon perspective suggests a hybrid
strategy: use deep models for raw predictive power, but use interpretable
trees (and their Rashomon sets) to understand \emph{which features matter
and why}.  The two approaches complement each other.
